% sec-A083835.tex -- k = 6 section of the unified kp+1
% paper. Standalone fragment, \input by
% A083833-UNIFIE-paper.tex.
% Environments and notation inherited from sec-carry.tex; the
% guards below are no-ops once defined there.
\providecommand{\IsPal}{\mathrm{IsPal}}
\makeatletter
\@ifundefined{theorem}{\newtheorem{theorem}{Theorem}}{}
\@ifundefined{corollary}{\newtheorem{corollary}[theorem]{Corollary}}{}
\@ifundefined{remark}{\newtheorem{remark}[theorem]{Remark}}{}
\makeatother

\section{Multiplier $k = 6$: A083835 and its image A083836}
\label{sec:k6}

OEIS A083835~\cite{oeis} (Murthy and Srikanth, May 2003;
extended by R.~Chandler) lists the palindromes $p$ such
that $q = 6p+1$ is again a palindrome:
$1$, $9$, $99$, $121$, $131$, $999$, $9999$, $12021$,
$12121$, $13031$, $13131$, \ldots; A083836 is the image
sequence $7$, $55$, $595$, $727$, $787$, $5995$,
$59995$, $72127$, \ldots
No conjecture is recorded on either entry. At this
multiplier the shared device yields not only the
necessary structure but the complete classification,
proved in both directions in the kernel.

\begin{theorem}[complete classification of A083835]
\label{thm:k6-iff}
Let $\IsPal(p,m)$ with $m \ge 1$. Then $6p+1$ is a
palindrome of some length $m'$ if and only if exactly
one of the following holds:
\begin{itemize}
\item[(i)] $m' = m + 1$ and every digit of $p$ equals
$9$, i.e.\ $p = 99\ldots9$ is the repdigit $9\,R_m$;
\item[(ii)] $m' = m$, $m$ is odd, and the digits of $p$
obey the alternating classes $a_j \in \{0,1\}$ at even
positions and $a_j \in \{2,3\}$ at odd positions.
\end{itemize}
\end{theorem}

\begin{corollary}[image companions: A083836]
\label{cor:k6-img}
Under the same hypotheses, every digit of $q = 6p+1$
avoids $\{0,3,4,6\}$, i.e.\ lies in $\{1,2,5,7,8,9\}$.
If $m'$ is even then $q$ has the exact form
$5\,9^{\,m'-2}\,5$: endpoints $5$, all interior digits
$9$ (the images of the repdigit branch).
\end{corollary}

\begin{proof}
Necessity is the instantiation of
Section~\ref{sec:carry} at $k = 6$: carries
$c_j = \bigl\lfloor\bigl(6\,(p \bmod 10^j)+1\bigr)
/10^j\bigr\rfloor \le 5$ with $c_0 = 1$. At constant
length, Lemma~\ref{lem:same} gives $m$ odd, $a_0 = 1$
and the classes $\{0,1\}$/$\{2,3\}$; at overflow, the
head equation $(6a_0+1) \bmod 10 = c_m \in \{1,\dots,5\}$
has the single solution $a_0 = 9$, $c_m = 5$, and the
mirrored local system \emph{(A)}+\emph{(B)} of
Lemma~\ref{lem:overflow} has the unique fixed point
$(a',c') = (9,5)$ -- verified by a finite check over the
ten digits, with no residual spurious state (unlike
$k = 3$): every digit is $9$ and
$q = 5\,9^{\,m-1}\,5$. For the converse directions the
carry chain is rebuilt from the prescribed digits: an
all-$9$ palindrome of length $m$ gives
$6p+1 = 5\,9^{\,m-1}\,5$, a palindrome of length $m+1$;
an odd-length palindrome with the alternating classes
closes the chain $c_j = 1$ at even $j$, $c_j = 0$ at odd
$j$, gives $c_m = 0$ and a palindromic $6p+1$ of length
$m$. Both directions are kernel-verified
(Table~\ref{tab:lean}). The image alphabet reads off
the branches: same-length digits in $\{1,2,7,8\}$,
overflow digits in $\{5,9\}$.
\end{proof}

\begin{remark}[the even-length members of A083835]
\label{rem:k6-even}
Since the same-length branch requires $m$ odd, every
even-length member of A083835 is the repdigit
$99\ldots9$, with image $5\,9^{\,m'-2}\,5$; the
published even-length terms are exactly $99$, $9999$,
$999999$, $99999999$. The same parity argument gives
$q$ even-length only in the overflow branch, which is
the second clause of Corollary~\ref{cor:k6-img}.
\end{remark}

\begin{remark}[the shared alternating classes at
$k = 5$ and $k = 6$]
\label{rem:converse}
Table~\ref{tab:classes} shows the same-length classes
coincide at $k = 5$ and $k = 6$: $\{0,1\}$ at even
positions, $\{2,3\}$ at odd positions, $a_0 = 1$. Hence
every non-repdigit member of A083833 carries a digit
structure which Theorem~\ref{thm:k6-iff} certifies to
be a member of A083835: the inclusion
\[
\{\text{non-repdigit terms of A083833}\}
\;\subseteq\;
\{\text{non-repdigit terms of A083835}\}
\]
is kernel-verified, and the $29$ published terms of the
two entries display the same non-repdigit members
($121, 131, 12021, 12121, 13031, \ldots$). The reverse
inclusion follows on paper: for the same classes the
$k = 5$ chain closes with the same alternating carries
$1,0,1,0,\ldots$ (the generic instance is
Lemma~\ref{lem:altsuff}) -- the exact analogue of the certified
$k = 6$ converse, and likewise kernel-certified at
$k = 5$ by \texttt{alt\_gives\_pal}
(Remark~\ref{rem:sufftask}).
\end{remark}
